% paper_edits_2026-11-04.tex: edit plan for V_Friction.tex. Pedro applies each block by hand: find the SEARCH text, replace it with the REPLACE text.
% BASIS: your save of 11-08 00:32 (X1 to X3, E3a, E3b through the advertising exclusion applied, with slips). Labels follow your numbering:
%          eq:12 to eq:21 are in the paper; E3c adds eq:22 (Pix exposure, openings). Time-deposit parts are 4^{m} (matured) and 4^{d} (daily liquidity).
% PENDING: Y1, Y2, Y3, Z1, Z2, Y4, E3c, E4, in that order (listed in that order below); then the bibliography at the end.
%          Assumptions are an enumerated list A1 to A9 in Appendix B (labels assum:*), referred to with \ref as the estimation methods are.
% TEST (11-08 00:37, sandbox with a biblatex stub; macros of 11-08 00:20): blocks unique, every \pn key resolves except the two branch-presence keys requested from IMPR, no undefined reference, no overfull box
%          from the edits. After applying: compile with biber, then python make_paper_equations.py.
% NOTATION: time-deposit parts 4^{m}, 4^{d}; age a_i, horizon H, clock premium Delta^c(a_i,H_i); scales varsigma, varsigma_varnothing (code: a, a_null); effects psi;
%          exits vartheta; one kappa-tilde per source pair (2 savings, 1 checking, 4 daily-liquidity time deposits); the depositor's source (j'(i),k'(i)).

% ============================================================================================================================
% Apply in this order: Y1, Y2, Y3, Z1, Z2, Y4, E3c, E4. Each SEARCH exists in your 00:32 save. Then add the four bibliography
%     entries at the end of this file to References.bib.
% ============================================================================================================================

% ========================================================================================================================
% Y1. Section 7.2, Pix, the mandate sentence. required. 'or' should be 'of'.
% ========================================================================================================================
% SEARCH
so a comparison or source pools on and off Pix within a market is empty.
% REPLACE
so a comparison of source pools on and off Pix within a market is empty.

% ========================================================================================================================
% Y2. Section 7.2, Pix, the footnote. required. The year (2002Q4) and the key name (rosted) are mistyped; the key prints ?? as is.
% ========================================================================================================================
% SEARCH
In 2002Q4 the \pn[w]{pix.rosted.B.n_mandated}
% REPLACE
In 2020Q4 the \pn[w]{pix.roster.B.n_mandated}

% ========================================================================================================================
% Y3. Section 7.2, Pix, the footnote. required. 'form' should be 'from'.
% ========================================================================================================================
% SEARCH
in any quarter form 2020Q4 to 2024Q4.
% REPLACE
in any quarter from 2020Q4 to 2024Q4.

% ========================================================================================================================
% Z1. Section 7.2, Pix, after the equation for kappa-tilde_mt. required. Defines the launch term and the dose kappa-tilde^U.
% ========================================================================================================================
% SEARCH
with $\hat{U}_m$ the market's Pix use predicted by its 2019 mobile and 4G/5G access. The before-after change and the dose are two readings of the same object.
% REPLACE
with $\hat{U}_m$ the market's Pix use predicted by its 2019 mobile and 4G/5G access. The launch term $\tilde{\kappa}^{\mathrm{post}}-\tilde{\kappa}^{\mathrm{pre}}$ is the change common to every market, and the dose $\tilde{\kappa}^{U}$ is the further change per unit of predicted use. They are two readings of the same object.

% ========================================================================================================================
% Z2. Section 7.2, Pix, after the equation for alpha_t. required. Defines alpha^{PU}.
% ========================================================================================================================
% SEARCH
The shift $\alpha^{P}$ is a common change at the launch in the response of every pool to the price. It does not measure attention, which $\psi_{mt}$ absorbs.
% REPLACE
The shift $\alpha^{P}$ is a common change at the launch in the response of every pool to the price, and $\alpha^{PU}$ is its further change per unit of predicted use. Neither measures attention, which $\psi_{mt}$ absorbs.

% ========================================================================================================================
% Y4. Section 7.2, advertising. required. A missing article.
% ========================================================================================================================
% SEARCH
This exclusion is supported by  finding of
% REPLACE
This exclusion is supported by the finding of

% ========================================================================================================================
% E3c. Section 7.2 (rest of advertising, stakes, fees and sweeps) and 7.3. required. Replaces your note and the unfinished sentence
%     on the five largest conglomerates. 'Known everywhere' was the awareness note's identifying assumption A_b = 1, stated as a fact.
%     It is now stated as an assumption, with its support: branch presence (two keys requested from IMPR, MAIN_REPLY 11-08 00:35,
%     ?? until they exist) and Datafolha Top of Mind (DATA_RAW/Datafolha/top_of_mind: Banco do Brasil first in every edition
%     2016 to 2022; in 2020 the five are the five banks named first, 30, 18, 18, 17 and 8 percent).
% ========================================================================================================================
% SEARCH
works mainly through awareness. %NOTE TO PEDRO: CHECK THIS!!
The give largest conglomerates are known
% REPLACE
works mainly through awareness.

We assume that every saver is aware of the five largest conglomerates, Banco do Brasil, Bradesco, Caixa, Ita\'{u} and Santander. In 2019Q4, at least one of them has a branch in \pn[p0]{desc.big5.branch.any.mca_share.2019Q4} percent of markets, which hold \pn[p0]{desc.big5.branch.any.pop_share.2019Q4} percent of the population. In Datafolha's annual survey of spontaneous brand recall, Banco do Brasil is the bank named first in every edition from 2016 to 2022, and in 2020, the one edition that reports all five, they are the five banks named first \parencite{datafolha_topofmind}.

Branch presence and recall do not show that every saver knows all five, so the normalization is an identifying assumption rather than a measurement. If $\kappa$ and $h$ are the same for these five and for the other institutions, the gap between the other institutions' $\tilde{\kappa}$ and theirs reads $-\ln\bigl[\chi_{jmt}+A_{jmt}(1-\chi_{jmt})\bigr]$.

The register of client relationships (CCS) gives a second moment, since opening a relationship at $j$ needs awareness of $j$ and moves no money. With $\mathrm{Rel}_{jt}$ the stock of relationships at $j$, $\mathrm{Adults}_t$ the adult population and $\mathrm{Openings}_{jt}$ the new relationships,
\begin{gather}
	\frac{\mathrm{Openings}_{jt}}{\mathrm{Adults}_t-\mathrm{Rel}_{jt-1}}=(1-\phi_t)\,A_{jt}\,P(j\mid\text{aware}), \label{eq:22}
\end{gather}
so openings respond to advertising through $A$ and to spreads through $P(j\mid\text{aware})$, the choice probability of $j$ among depositors who consider it.

The advertising accounts of every institution the BCB supervises, released upon request, give a quarterly series for nearly every conglomerate in the estimation sample, digital entrants included (Appendix \ref{app:data:advertising}). The model of Section \ref{sec:model} is nested as the full-awareness limit, in which depositors are not always awake, but aware of every available option.

\subsubsection{The Stakes}\label{sec:disentangle:attention}\label{sec:disentangle:attention:stakes}

Attention that responds to the stakes of shopping predicts that sleepiness falls when the policy rate rises. The lagged Selic coefficient of Section \ref{sec:estimation:sleep} has that sign. It is significant under \ref{estimation:single_idx_time} but not under the preferred \ref{estimation:single_idx}, so the current estimates lean toward this channel without establishing it.

In \eqref{eq:20} the awake share sits in the market-quarter effect $\psi_{mt}$, which absorbs any response of attention to the stakes, so $\tau_1$ is not identified there. It comes from the sleepiness step, with the stakes of \eqref{eq:15} as one more market state in the index of \eqref{eq:8} and \eqref{eq:10}. What \eqref{eq:20} tests is that the stakes shift neither $\tilde{\kappa}$ nor $\alpha$.

The stakes are endogenous to the same shocks as the price, so they enter instrumented by the cost shifters of Section \ref{sec:empirical:sleep}, which move the dispersion of spreads, with the lagged market structure as controls. Identification leans on cross-market within-quarter variation, rather than on the Selic time series alone, which is confounded with the aggregate cycle.

\subsubsection{Fees, Sweeps, and the Within-Bank Margin}\label{sec:disentangle:switching}

Account fees are the explicit cost of keeping an account, and they fall mostly on incumbents. They make an incumbent account costly to keep and a digital account nearly free (Appendix \ref{app:netflow:withinbank}). Fees therefore push balances toward digital banks, so they cannot explain why balances drift toward those banks slowly. A switching cost that slows the drift would be a hassle or an information cost rather than a fee.

Three features of the data bear on the comparison of moves within and across institutions. Balances are quarter-end stocks, so only net flows are observed, which is why \eqref{eq:20} is written on the net flow. The within-bank margin is not frictionless either, since moving a time deposit before maturity restarts its tax clock (Appendix \ref{app:clock}) and a new product may require onboarding at the bank.

Automatic sweeps work the other way, since they move balances within a bank without a decision. Since 2022, Santander's general conditions for new individual accounts make its sweep from checking accounts into certificates of deposit a default feature \parencite{santander_condicoes_pf_2022}, while Bradesco's sweep requires the client to opt in \parencite{bradesco_investfacil_2015,bradesco_investfacil_2020}. The checking pools enter \eqref{eq:20} with their own scale $e^{\varsigma}$ for that reason, and the dates of the default sweeps are known. ESTBAN is also published by branch, monthly from 2013, so moves between a bank's own branches can be followed as well.

\subsection{Evidence So Far}\label{sec:disentangle:evidence}

Three facts of the net-flow panel set the scene. At the median time-deposit cell, the net flow is \pn[p0]{kappa.k0.moving.median.k4} percent of the carried stock, so the margin \eqref{eq:20} explains is not small.
The outside pool is \pn[p0]{kappa.k0.outside.share_Mhat.median} percent of the market size at the median market-quarter, so it is the largest pool in the market and cannot be left out of the equation.
And \pn[p0]{kappa.k0.logR2.between_share} percent of the variance of the log rival-savings weight is between bank $\times$ market cells, so the level of $\tilde{\kappa}$ rests on the small part of that variance within cells.

One discrepancy is recorded rather than resolved. In the world the estimates of Section \ref{sec:estimation} imply, an awake R\$ re-chooses its own product with probability \pn[2]{kappa.k2.world.stay_prob.median} and awake money is \pn[p1]{kappa.k2.world.awake_share} percent of deposits, against a net flow of \pn[p0]{kappa.k0.moving.median.k4} percent of the carried stock in the data.
The design is built on the net flow, which is observed. The gap suggests that the sleepy share is too high in many cells or that asleep balances also move, the two readings Section \ref{sec:empirical:demand} gives for the truncated active balances. It is revisited with the random-coefficients refit.

Appendix \ref{app:simulation} simulates \eqref{eq:17} on the panel's pools at the mean utilities, the sleepy shares and the market sizes of Section \ref{sec:estimation}. The switching cost is \pn[0]{kappa.k2b.truth.level} before 2020Q4 and changes by \pn[1]{kappa.k2b.truth.change} at the launch, the dose is \pn[2]{kappa.k2b.truth.dose} per unit of exposure, and the price coefficient is that of Section \ref{sec:estimation:demand}.

Table \ref{tab:k2b_pairs} reports what \eqref{eq:20} recovers, in a world where every pool has the same denominator and in one where the denominators differ as they do in the data. The headline object is recovered. For the savings pair, the change lies between \pn{kappa.k2b.sav.change.fe_ba} and \pn{kappa.k2b.sav.change.fe_dose} in the four fits, and the dose is \pn{kappa.k2b.sav.dose.fe_dose} and \pn{kappa.k2b.sav.dose.ideal_dose}.

The price coefficient is recovered (\pn{kappa.k2b.alpha.fe_ba} and \pn{kappa.k2b.alpha.ideal_ba}), and its shift at the launch, $\alpha^{P}$, is at most \pn{kappa.k2b.alpha_P.max} in absolute value.
The level is less secure. In three of the four fits it lies below the truth by more than its standard error, for instance at \pn{kappa.k2b.sav.level.fe_dose} (\pn{kappa.k2b.sav.level_se.fe_dose}) and \pn{kappa.k2b.sav.level.ideal_ba} (\pn{kappa.k2b.sav.level_se.ideal_ba}), and the dose fit with pool-specific denominators stopped at its iteration cap.

The checking pair is not pinned down even where the common denominator holds. The time-deposit pair recovers the level and the change, but its dose is \pn{kappa.k2b.td.dose.fe_dose} and \pn{kappa.k2b.td.dose.ideal_dose}, on opposite sides of the truth, so neither pair is a second estimate.

A linearized version of \eqref{eq:20} recovers nothing about $\kappa$ once the entry scale varies across cells, which is why the equation is estimated nonlinear in the price.

Three breaks of the assumptions bear on the level (Table \ref{tab:k2_breaks}). When tastes vary across depositors as in the random coefficients of Section \ref{sec:estimation:demand}, the level fails. For a true $\kappa$ of \pn[1]{kappa.k2.rc.k05.truth} and \pn[0]{kappa.k2.rc.k10.truth}, a single $\tilde{\kappa}$ shared by the pairs is \pn{kappa.k2.rc.k05.rc.kappa} (\pn{kappa.k2.rc.k05.rc.se}) and \pn{kappa.k2.rc.k10.rc.kappa} (\pn{kappa.k2.rc.k10.rc.se}), and neither fit converges, while the same rows with logit tastes give \pn{kappa.k2.rc.k05.logit.kappa} and \pn{kappa.k2.rc.k10.logit.kappa}. The random coefficients move the entry scale within cells, which breaks \ref{assum:scale}.

Sorting toward the incumbent raises $\tilde{\kappa}$ above $\kappa$ in every simulated case, and by more than $-\ln h$ once a positive cost and sorting combine. Leaving out the outside pool drives $\tilde{\kappa}$ to about zero whatever the true $\kappa$. None of these breaks was simulated for the change at the launch, which rests on \ref{assum:launch}.

On the opening margin, advertising moves openings at the $\mathrm{B}$ institutions outside the five largest, \pn{pix.adv.rest.elast} per log unit in the quarter (bootstrap $p$ \pn{pix.adv.rest.elast_p}) and \pn{pix.adv.rest.sum_lags} over three quarters. At the five largest the elasticity is \pn{pix.adv.big5.elast}, with a bootstrap $p$ of \pn[2]{pix.adv.big5.elast_p}. That is consistent with full awareness of the largest institutions, although the difference between the two groups has not been tested.

Instrumented by its own four-quarter lag, the elasticity is \pn{pix.adv.iv.elast} (first-stage $F$ \pn[1]{pix.adv.iv.fs_f}, Anderson and Rubin confidence set \pn{pix.adv.iv.ar.lo} to \pn{pix.adv.iv.ar.hi}), so same-quarter simultaneity does not overstate it, and the advertising of the other institutions in a conglomerate's group lowers the institution's own openings (\pn{pix.adv.group_placebo.coef}), as rivals' awareness should.
Branches move too little within a cell to identify $\gamma_2$, so awareness rests on advertising.

The stakes of \eqref{eq:15} have a median of R\$\pn[0]{pix.stakes.median} a year, and the cost shifters predict their log with an $F$ of \pn{pix.stakes.fs.f}, or \pn{pix.stakes.fs.f_net} net of the lagged market structure. Part of that strength comes from the lagged deposit weights that any bank attribute would carry: shuffling the cost-predicted spreads across institutions within the quarter gives a median $F$ of \pn{pix.stakes.fs.f_placebo}, and \pn[p0]{pix.stakes.fs.placebo_share_above} percent of the shuffles reach the real statistic, so the cost shifters carry the rest and the instrument for the stakes exists.

The switching-cost equation itself runs next, by source pair and with the Pix and advertising shifters. Pilots on a sample of markets do not yet pin down the change at the launch: the objective, profiled over the effects, is rough, and the estimate depends on its starting value. The full estimation waits for a smoother solve of the effects.

The time-deposit pair and the clock wait for the registry's split of time deposits by liquidity, requested from the BCB. The random-coefficients refit with the source follows once \eqref{eq:20} has said which frictions to carry.

% ========================================================================================================================
% E4. End of file. required. Appendices A (the tax clock), B (the net-flow equation) and C (simulation evidence), before
%     \end{document}. Appendix B now has subsections: B.1 derivation, B.2 the scales and the exit, B.3 identification with two
%     propositions and proofs, B.4 estimation and inference, B.5 the Pix exposures with Table tab:pix_exposure_fs, B.6 the within-bank
%     margin. Cites rothenberg1971identification (entry at the end of this file).
% ========================================================================================================================
% SEARCH
\end{document}
% REPLACE
\section{The Tax Clock}\label{app:clock}

Section \ref{sec:inst_setting:deposits} says that moving a time deposit before maturity restarts the tax clock. This appendix prices the clock from the law and says how it enters the model of Section \ref{sec:disentangle:model} and the equation of Section \ref{sec:disentangle:netflow}. Let $w(n)$ be the withholding rate on income at age $n$ days, the regressive schedule of \textcite{brasil_lei11033_2004},
\begin{gather}
	w(n)=\begin{cases}0.225 & n\le 180\\ 0.20 & 180<n\le 360\\ 0.175 & 360<n\le 720\\ 0.15 & n>720,\end{cases} \label{eq:app_tau}
\end{gather}
and $\iota(n)$ the share of income taken by the tax on financial operations (IOF), 0.96 on day 1 falling to 0 on day 30 \parencite{brasil_decreto6306_2007}. Let $r$ be the application's annual gross rate, the institution's own time-deposit rate (in this appendix only, since elsewhere $r_t^f$ is the risk-free rate), $a$ the application's age in days when the decision is taken, and $H$ the horizon in days the money then stays invested. Staying and redeeming at age $a+H$ leaves the depositor with
\begin{gather}
	\mathcal{W}^{\mathrm{stay}}(a,H)=1+\bigl[1-w(a+H)\bigr]\Bigl[(1+r)^{(a+H)/365}-1\Bigr] \label{eq:app_Ws}
\end{gather}
per R\$ of principal. Moving now to an institution paying $r^\prime$ means redeeming, paying the clock, reinvesting the net proceeds $V(a)$, and redeeming at the new application's age $H$,
\begin{gather}
	V(a)=1+\bigl[1-w(a)\bigr]\bigl[1-\iota(a)\bigr]\Bigl[(1+r)^{a/365}-1\Bigr], \label{eq:app_P} \\
	\mathcal{W}^{\mathrm{move}}(a,H;r^\prime)=V(a)\Bigl\{1+\bigl[1-w(H)\bigr]\Bigl[(1+r^\prime)^{H/365}-1\Bigr]\Bigr\}. \label{eq:app_Wm}
\end{gather}
The cost of the clock, evaluated at $r^\prime=r$, the break-even rate $r^{*}$ at which moving pays, and the clock premium $\Delta^{c}$ are
\begin{gather}
	C(a,H)=\mathcal{W}^{\mathrm{stay}}(a,H)-\mathcal{W}^{\mathrm{move}}(a,H;r), \nonumber \\
	\mathcal{W}^{\mathrm{move}}(a,H;r^{*})=\mathcal{W}^{\mathrm{stay}}(a,H), \qquad \Delta^{c}(a,H)=r^{*}-r. \label{eq:app_cost}
\end{gather}
Table \ref{tab:tax_clock} computes both for a deposit paying the CDI rate, moved to an institution paying the same, at the level of the CDI in its Panel A. The cost runs from \pn[2]{clock.t3.cost.lo} to \pn[2]{clock.t3.cost.hi} R\$ per R\$100 and the premium from \pn[2]{clock.t3.premium.lo} to \pn[2]{clock.t3.premium.hi} percentage points a year, or \pn[1]{clock.t3.premium_cdi.lo} to \pn[1]{clock.t3.premium_cdi.hi} percent of the CDI, over the cells of age and horizon in that panel.

A deposit paying less than the CDI behaves like a lower level of the CDI, with a smaller cost and premium in points but a similar premium as a percent of its rate. Ages and horizons are illustrative cells, not the distribution of the stock, which the registry request to the BCB would give.

\begin{table}[H]
	\centering
	\begin{threeparttable}
		\caption{The Cost of Moving a Time Deposit Before Maturity}\label{tab:tax_clock}
		\input{tab_tax_clock.tex}
		\begin{tablenotes}[flushleft]
			\item \footnotesize \textit{Notes:} A daily-liquidity time deposit paying the CDI is redeemed at the age shown and reinvested at another bank paying the same rate for the horizon shown. Cost: after-tax wealth lost per R\$100 of principal relative to staying. Income is taxed at 22.5 percent up to 180 days, 20 percent up to 360, 17.5 percent up to 720 and 15 percent beyond \parencite{brasil_lei11033_2004}, and redemptions within 30 days also pay the tax on financial operations (IOF) \parencite{brasil_decreto6306_2007}, which does not bind at these ages. Break-even premium: the rate above the CDI the new bank must pay over the horizon for moving to leave the depositor where staying would, in percentage points a year and as a percent of the CDI. Panel B gives the range over the cells of Panel A at three levels of the CDI. A deposit paying less than the CDI behaves like a lower level of the CDI, with a smaller cost and premium in points but a similar premium as a percent of its rate. Ages and horizons are illustrative, not the distribution of the stock. A time deposit cannot renew automatically \parencite{cmn_resolucao_3454_2007,cmn_resolucao_5005_2022}, so at maturity the clock restarts whether the money stays or moves, and the cost is zero on that date.
		\end{tablenotes}
	\end{threeparttable}
\end{table}

\paragraph{In the model.} The clock is a cost on leaving a time deposit before maturity for another inside product. It enters the utility of every inside option other than staying, and not the outside option, since redeeming to spend restarts no clock. In the units of $\delta$ the cost is $-\bar{\alpha}\,\Delta^{c}_{it}$, with $\Delta^{c}_{it}\equiv\Delta^{c}(a_i,H_i)$ at the balance's age and horizon. It is the premium that makes moving pay times the utility per percentage point, a positive number as in Section \ref{sec:disentangle:model}. That is the second line of \eqref{eq:12}.

The premium is derived for a move into another time deposit. A move into savings, whose income is untaxed \parencite{brasil_lei8981_1995}, would need its own premium, and the model uses one premium for both destination types as an approximation.

Of the type-4 stock an institution carries into a quarter, three parts face different costs. A share $\omega_{jt}$ matures in the quarter and lands in the checking account whether its holder is awake or asleep, since a time deposit cannot renew automatically. The National Monetary Council (CMN) banned automatic renewal in 2007 \parencite{cmn_resolucao_3454_2007} and kept the ban in the resolution that replaced that rule on 2 May 2022 \parencite{cmn_resolucao_5005_2022}. Matured money pays nothing to move. A share $(1-\omega_{jt})\lambda_{jt}$ has daily liquidity and pays $C(a,H)$ to move. The rest, $\ell_{jt}=(1-\omega_{jt})(1-\lambda_{jt})$, is locked and cannot move even when awake, which is the lock in the carry term of \eqref{eq:17}.

\paragraph{In the net-flow equation.} The equation of Section \ref{sec:disentangle:netflow} never uses $w(n)$ or $\iota(n)$. It splits the time-deposit pools by their composition, $\omega_{jt}$ and $\lambda_{jt}$, which are data, and reads the clock from the gap between source pairs. Staying is exempt, so own daily-liquidity money that stays enters at weight one. Rivals' daily-liquidity money pays the clock, so its rival factor in \eqref{eq:20} is $e^{-\tilde{\kappa}_4}=e^{-\kappa}E_c$, with
\begin{gather}
	E_c\equiv\mathbb{E}\bigl[e^{\bar{\alpha}\,\Delta^{c}(a,H)}\bigr], \label{eq:app_clockfactor}
\end{gather}
where the expectation is over the ages and horizons of the stock and the exponential is what the cost does inside a logit. Under \ref{assum:one_kappa}, and with awareness and sorting common to the two pairs, $\ln\hat{E}_c=\tilde{\kappa}_2-\tilde{\kappa}_4$.

The law gives a test. The estimate should equal \eqref{eq:app_clockfactor} over the distribution of ages and horizons in the stock, which exceeds $e^{\bar{\alpha}\,\mathbb{E}[\Delta^{c}(a,H)]}$ by Jensen's inequality. The comparison is therefore with the expectation of the exponential, not with the exponential at the mean premium. Savings money carries no cost in the model, so the gap between the savings and the time-deposit pairs measures the clock alone.

The exit takes the movable stock with one market-quarter coefficient (Appendix \ref{app:netflow:terms}). The share $\omega_{jt}$ comes from the statements' bucket of maturities within three months, annual for the institutions that publish it and imputed by firm type for the rest. The statements' line for deposits without maturity puts $\lambda_{jt}$ near zero for every $\mathrm{B}$ institution, probably because time deposits with daily liquidity carry contractual maturities. The headline therefore sets $\lambda_{jt}=0$, where neither $\tilde{\kappa}_4$ nor $E_c$ is estimable, and a robustness row sets $\lambda_{jt}=1$. A registry request to the BCB asks for $\lambda_{jt}$ directly and for early against at-maturity redemptions, which would give the exit flow split by clock status.

\section{The Net-Flow Equation}\label{app:netflow}

This appendix derives \eqref{eq:20} from \eqref{eq:18}, names each assumption, and says what each buys and what it costs if false. It then sets out the scales and the exit coefficient, the identification of $\tilde{\kappa}$, the estimator, the Pix exposures and the within-bank margin.

Throughout, a pool is a source product $p=(j^\prime,k^\prime)$ in a market, with time deposits split into their matured part $4^{\mathrm{m}}$ and their daily-liquidity part $4^{\mathrm{d}}$, and $D_p$ is the denominator of its conditional share.

The estimating equation rests on nine assumptions, referred to by their labels throughout.
\begin{enumerate}[label = A\arabic*, ref = A\arabic*]
	\item\label{assum:mean_utility} Mean utility: the conditional shares are evaluated at the mean utilities, $\mu_{ijkmt}=0$.
	\item\label{assum:common_denominator} Common denominator: every pool in a market has the denominator $D_{mt}$.
	\item\label{assum:source_cost} Source cost: no source type carries a cost of moving other than the clock.
	\item\label{assum:one_kappa} One switching cost: each source pair carries one $\kappa$ for the market, stated in $e^{-\kappa}$.
	\item\label{assum:market_size} Market size: the market size of Appendix \ref{app:estimation:blp:marketsize} measures the outside pool.
	\item\label{assum:scale} Common scale: an institution's entry scale is the same for every pool that feeds it and does not move within a bank $\times$ market cell beyond the market-quarter effect.
	\item\label{assum:awareness} Awareness: savers at rival institutions are aware of $j$, $A_{jmt}=1$.
	\item\label{assum:sorting} Sorting: tastes are sorted toward the incumbent only, $h_{jmt}\in(0,1]$.
	\item\label{assum:launch} Launch: awareness and sorting do not move at the launch of Pix in 2020Q4.
\end{enumerate}

\subsection{Derivation}\label{app:netflow:derivation}

\paragraph{\ref{assum:mean_utility}, evaluation at the mean utility.} Set $\mu_{ijkmt}=0$ in \eqref{eq:13}, so the conditional share is the plain logit of \textcite{berry1994estimating}. The exact object is an integral over the tastes of the holders of pool $p$,
\begin{gather}
	P(j4\mid p)=\int\frac{\exp\bigl\{v_{ij4mt}\bigr\}}{D_{p}(\mu_i)}\,\text{d}F_{p}(d_{imt},\nu_{imt}), \label{eq:app_Pexact}
\end{gather}
with $v_{ij4mt}$ the deterministic utility of \eqref{eq:12} at the pool's source, $F_p$ the holders' distribution of demographics and draws, and $D_p(\mu_i)$ the pool's denominator at the depositor's tastes. The share \eqref{eq:13} integrates over the market's $F$ instead, and the difference is the sorting defined below.

What \ref{assum:mean_utility} buys is an equation that needs no $\theta_2$. Without it, \eqref{eq:app_Pexact} is a function of $\theta_2$, and the estimates of Section \ref{sec:estimation:demand} cannot stand in, since they were estimated at $\kappa=0$. Dropping \ref{assum:mean_utility} means refitting the share system with $\kappa$. Appendix \ref{app:simulation} measures the cost of \ref{assum:mean_utility} with the random coefficients of Section \ref{sec:estimation:demand} in the data-generating process, and the level of $\tilde{\kappa}$ does not survive it.

\paragraph{\ref{assum:common_denominator}, a common denominator.} Under \ref{assum:mean_utility}, $D_p=1+\sum_{(l,k)}\exp\{v_{ilkmt}\}$ differs across pools through which institution's products carry no $\kappa$ and, for a daily-liquidity time-deposit pool, which products carry the clock. \ref{assum:common_denominator} says that difference is negligible, $D_p\approx D_{mt}$ for every pool in the market. It holds when the outside option dominates, $D_p\approx 1+\sum_{(l,k)}\exp\{\delta_{lkmt}\}$ with the inside terms small, so every pool faces the same menu to first order.

Whether it holds is an empirical fact about the menu. On the estimates of Section \ref{sec:estimation:demand}, the outside option is \pn[p0]{kappa.k0.outside.menu_share.median} percent of the median market's denominator, and an awake R\$ re-chooses its own product with probability \pn[2]{kappa.k2.world.stay_prob.median}, so the inside terms are small. Under \ref{assum:common_denominator},
\begin{gather}
	P(j4\mid p)\approx c_{p}\,\frac{e^{\delta_{j4mt}}}{D_{mt}}\,e^{-\kappa\,\indicator{j^\prime\neq j}}, \label{eq:app_logit}
\end{gather}
where $c_p$ is the clock factor, $E_c$ of \eqref{eq:app_clockfactor} for a daily-liquidity time-deposit pool that moves and one for every other pool and for money that stays.

What \ref{assum:common_denominator} buys is separability. The entry from a pool is the pool's size times $e^{-\kappa\indicator{j^\prime\neq j}}$ times a function of the destination alone, which makes the rival factor read as $e^{-\kappa}$ and the equation estimable with fixed effects. If \ref{assum:common_denominator} is false, the rival-to-own ratio is $e^{-\kappa}D_{j}/D_{j^\prime}$, so the error in $\kappa$ is $\ln(D_{j^\prime}/D_{j})$.

In the simulation the ratio implied by the simulated denominators sits slightly above $e^{-\kappa}$ for the savings pair, so \ref{assum:common_denominator}'s error pushes $\tilde{\kappa}$ down. The error differs across pools, which is one reason each source pair has its own $\tilde{\kappa}$. It is not the only reason the pairs differ, since the checking pair is weakly identified even where \ref{assum:common_denominator} holds (Appendix \ref{app:simulation}).

Awake daily-liquidity money at $(j,4)$ stays with probability $e^{\delta_{j4mt}}/D_{mt}$, since staying pays no cost, and otherwise leaves. In the world the estimates imply that probability is small, so the exit is nearly the whole awake daily-liquidity stock, and the stayers re-enter through the own daily-liquidity pool on the entry side. Relative to own savings their weight is one, since the savings forfeit cancels (Section \ref{sec:disentangle:model}). Matured money leaves whether its holder is awake or asleep, and its awake part re-enters as checking money.

\paragraph{\ref{assum:source_cost}, no source type carries a cost other than the clock.} The savings forfeit falls on every use of the money, the outside option included, so with a common denominator it cancels (Section \ref{sec:disentangle:model}). Checking costs nothing to leave. The clock on a daily-liquidity time deposit is the same tax whichever institution receives the money, and staying is exempt. So
\begin{gather}
	P(j4\mid j,k^\prime)=c_{p}\,\frac{e^{\delta_{j4mt}}}{D_{mt}},\qquad
	P(j4\mid j^\prime,k^\prime)=e^{-\kappa}\,c_{p}\,\frac{e^{\delta_{j4mt}}}{D_{mt}},\qquad j^\prime\neq j, \label{eq:app_Pown}
\end{gather}
with $p=(j,k^\prime)$ and $p=(j^\prime,k^\prime)$.

In the model every source has the same scale. The estimating equation frees two scales against savings money as nuisance parameters, $\varsigma$ for checking money and matured time deposits, which have landed in checking, and $\varsigma_{\varnothing}$ for new money. Appendix \ref{app:netflow:terms} gives the reasons and what each scale absorbs. Own daily-liquidity time deposits that stay enter at weight one. For rivals' daily-liquidity time deposits the clock factor $E_c$ joins $e^{-\kappa}$ in $e^{-\tilde{\kappa}_4}$.

\paragraph{The pools and the price.} The destination is an institution's time deposits, and the sources are every pool in the market. Own savings enters at weight one. Own checking enters at weight $Q^{1}$ and scale $e^{\varsigma}$, and a default sweep moves it into the own time deposit without a decision. Own matured time deposits have landed in checking, so they enter at weight $Q^{4^{\mathrm{m}}}$ and the same scale. Own daily-liquidity time deposits re-enter at weight $Q^{4^{\mathrm{d}}}$ and scale one when they stay, and they also make the exit.

Rivals' savings enter at weight $R^{2}$ with factor $e^{-\tilde{\kappa}_2}$. Rivals' checking and matured time deposits enter at weights $R^{1}$ and $R^{4^{\mathrm{m}}}$, with scale $e^{\varsigma}$ and factor $e^{-\tilde{\kappa}_1}$. Rivals' daily-liquidity time deposits enter at weight $R^{4^{\mathrm{d}}}$ with factor $e^{-\tilde{\kappa}_4}$, which carries the clock. The outside pool has no incumbent and pays no $\kappa$, and it enters at weight $\mathrm{Out}$ and scale $e^{\varsigma_{\varnothing}}$.

Each pair carries one $\kappa$ for the market (\ref{assum:one_kappa}), stated in $e^{-\kappa}$. If the cost differs across pairs of institutions, $e^{-\tilde{\kappa}}$ is the pool-weighted mean of $e^{-\kappa}$, so $\tilde{\kappa}$ is at most the pool-weighted mean of $\kappa$ by Jensen's inequality.

The outside pool is the market size of Appendix \ref{app:estimation:blp:marketsize} less last quarter's accrued deposits (\ref{assum:market_size}). New money has no incumbent and pays no $\kappa$. Its weight loads on the own and the rival pools alike, so leaving it out biases $\tilde{\kappa}$ down, and it is not small, at \pn[p0]{kappa.k0.outside.share_Mhat.median} percent of the market size at the median market-quarter.

The equation covers $\mathrm{B}$ institutions only, because the pools are local. $\mathrm{D}$ institutions hold one national stock, so $\mathrm{D}$ sources drop out of the rival pools. The omitted entry is divided by each destination's own savings, so the effects do not absorb it, and we do not sign its bias. $\mathrm{D}$ destinations would need their own equation at the national level, which we have not estimated.

Prepaid balances are not among the source pools either. Institutions set their rate, but in the panel every $\mathrm{B}$ institution carries the CDI rate on them, so they add no variation in price, and they are \pn[1]{slide.types.prepaid.share} percent of balances.

With a common denominator, the only product-specific object in \eqref{eq:app_logit} is the destination's mean utility. The characteristics $x_j$ are slow and go to the bank $\times$ market effects. What varies across institutions within a market-quarter is the institution's own time-deposit spread and its unobserved quality, so the price is the gap $g_{jmt}$ of Section \ref{sec:disentangle:netflow}, and $e^{-\bar{\alpha}g}$ is the price part of $e^{\delta_{j4}}$ up to a market-quarter constant.

\paragraph{From the identity to the estimator.} Substitute \eqref{eq:app_Pown} into \eqref{eq:18}, divide by the own savings pool, and write the weights \eqref{eq:19}, with the accrual factors of the entry pools absorbed into the scales. The entry is every pool in the market, own pools at their source-type scales and rivals' pools also at $e^{-\kappa}$, times one destination-specific factor, the share of the menu that $j$'s time deposit takes. The exit is all of the matured money and the awake daily-liquidity money of $j$'s time deposits.

Three things then turn it into \eqref{eq:20}. The entry scale $(1-\phi_{mt})e^{x_j^\prime\bar{\beta}+\xi_{j4mt}}/D_{mt}$ becomes $e^{\psi_{jm}+\psi_{mt}}$. This is \ref{assum:scale}, which says that an institution's scale is the same for every pool that feeds it and does not move within a bank $\times$ market cell beyond the market-quarter effect.

The exit coefficients, one on the matured stock and $1-\phi_{mt}$ on the daily-liquidity stock, become one market-quarter coefficient $e^{\vartheta_{mt}}$ on the movable stock, because a fit that holds the matured coefficient at one does not converge (Appendix \ref{app:netflow:terms}). And each pair's factor is written $e^{-\tilde{\kappa}}$, which equals $e^{-\kappa}$ only if savers at rival institutions are aware of $j$ (\ref{assum:awareness}) and their tastes are not sorted toward the incumbent.

\subsection{The Scales and the Exit}\label{app:netflow:terms}

In the model, each R\$ responds to $j$'s rate in the same way whatever its source, apart from $\kappa$ and the clock. The savings forfeit falls on every use of the money and cancels (Section \ref{sec:disentangle:model}), and checking balances and new money carry no cost of leaving, so every source has scale one. Equation \eqref{eq:20} keeps savings money as the reference, at scale one, and frees the scale of checking money, $e^{\varsigma}$, and the scale of new money, $e^{\varsigma_{\varnothing}}$.

Checking balances also serve payments, so less of each R\$ in checking may be free to move into a time deposit in a given quarter. Sweeps move checking balances into time deposits within an institution without a decision (Section \ref{sec:disentangle:switching}). Matured time deposits are credited to checking (Section \ref{sec:inst_setting:deposits}), so they take the checking scale.

New money is the market size net of the accrued stocks, so its measure carries the error of the estimated market size $\widehat{M}_{mt}$, and $\varsigma_{\varnothing}$ takes up part of that error. In the rivals' checking pools, $\varsigma$ enters together with $\tilde{\kappa}_1$, as $e^{\varsigma-\tilde{\kappa}_1}$. Only the own checking pool, which carries $e^{\varsigma}$ alone, separates the two (Proposition \ref{prop:app_ident}), which is one reason the checking pair is less precise than the savings pair.

In the model, the matured part of $j$'s time deposits leaves in full, whether its holder is awake or asleep, and the daily-liquidity part leaves only when its holder is awake. The exit of \eqref{eq:18} thus has coefficient one on $Q^{4^{\mathrm{m}}}$ and $1-\phi_{mt}$ on $Q^{4^{\mathrm{d}}}$. Equation \eqref{eq:20} replaces both with one coefficient per market and quarter, $e^{\vartheta_{mt}}$, on the movable stock $Q^{4^{\mathrm{m}}}+Q^{4^{\mathrm{d}}}$.

This is a response to the data, not to the model. With the statements' $\lambda_{jt}$ near zero, the movable stock is the matured part, and a fit that holds its coefficient at one does not converge. It stalls, or the parameters run to their bounds.

We have not isolated the cause, and three features of the data are candidates. The matured share $\omega_{jt}$ comes from an annual maturity bucket in the financial statements and is imputed where the statements do not report it, so the quarterly matured stock is measured with error. A $\lambda_{jt}$ near zero may count as matured the outflows of time deposits with daily liquidity, which move with attention. And a sweep can move matured money, once credited to checking, back into the same institution's time deposits within the quarter, so less of it leaves on net.

The free coefficient absorbs all three, and the model's constant is reported as a check. The exit then no longer tests the model. It is still separated from the market-quarter effect $\psi_{mt}$, since the entry scales with every pool in the market and with the price, and the exit only with $j$'s own movable stock.

The statements put $\lambda_{jt}$ near zero for every $\mathrm{B}$ institution, so the headline sets $Q^{4^{\mathrm{d}}}=R^{4^{\mathrm{d}}}=0$, and the time-deposit pair and the clock are not estimable there. That reading is probably an artifact of the statements' line for deposits without maturity, since time deposits with daily liquidity carry contractual maturities. A robustness row sets $\lambda_{jt}=1$ for every institution and estimates $\tilde{\kappa}_4$ and $E_c$.

\subsection{Identification}\label{app:netflow:ident}

Write the conditional mean of \eqref{eq:20} as
\begin{gather}
	m_{jmt}(\theta)=E_{jmt}B_{jmt}-\tilde{M}_{jmt}\,e^{\vartheta_{mt}},\qquad E_{jmt}=e^{\psi_{jm}+\psi_{mt}-\alpha g_{jmt}},\qquad B_{jmt}=1+\sum_{k=1}^{6}b_{k}X^{k}_{jmt}, \label{eq:app_mean}
\end{gather}
where
\begin{gather*}
	X_{jmt}=\bigl(Q^{4^{\mathrm{d}}},\,Q^{1}+Q^{4^{\mathrm{m}}},\,R^{2},\,R^{1}+R^{4^{\mathrm{m}}},\,R^{4^{\mathrm{d}}},\,\mathrm{Out}\bigr),\qquad b=\bigl(1,\,e^{\varsigma},\,e^{-\tilde{\kappa}_2},\,e^{\varsigma-\tilde{\kappa}_1},\,e^{-\tilde{\kappa}_4},\,e^{\varsigma_{\varnothing}}\bigr)
\end{gather*}
are the pool weights and their coefficients. The movable stock at its accrued value is $\tilde{M}_{jmt}=\bigl(1+r^{\mathrm{dep}}_{j4mt-1}\bigr)\bigl(Q^{4^{\mathrm{m}}}+Q^{4^{\mathrm{d}}}\bigr)$.

Own savings enters $B_{jmt}$ with coefficient one, which sets the scale of the effects. The first coefficient of $b$ is one as well, since own daily-liquidity money that stays pays no cost. The control-function residual of the spread enters the exponent of $E_{jmt}$, and its coefficient is suppressed from the notation. The parameters are $\theta=(\alpha,\ln b_2,\dots,\ln b_6,\psi,\vartheta)$.

\begin{prop}[Identification of the rival factors]\label{prop:app_ident}
	Let \ref{assum:mean_utility} to \ref{assum:scale} hold, and let the error of \eqref{eq:20} have mean zero given the pool weights, the gap, the movable stock and the control-function residual. Suppose that, at the true $\theta$, the variables $X^{k}_{jmt}/B_{jmt}$ for $k=2,\dots,6$, the gap $g_{jmt}$, and the ratio $\tilde{M}_{jmt}/(E_{jmt}B_{jmt})$ interacted with market $\times$ quarter indicators are linearly independent after their bank $\times$ market and market $\times$ quarter means are removed. Then $\theta$ is locally identified, up to one normalization of the effects in each market, and
	\begin{gather*}
		\varsigma=\ln b_2,\qquad \tilde{\kappa}_2=-\ln b_3,\qquad \tilde{\kappa}_1=\ln b_2-\ln b_4,\qquad \tilde{\kappa}_4=-\ln b_5,\qquad \varsigma_{\varnothing}=\ln b_6.
	\end{gather*}
	Where $Q^{4^{\mathrm{d}}}=R^{4^{\mathrm{d}}}=0$, as in the headline row, $b_5$ drops out and $\tilde{\kappa}_4$ is not identified.
\end{prop}

\begin{proof}
	Local identification holds if the Jacobian of $m(\theta)$ has full column rank \parencite{rothenberg1971identification}. Its columns are $E_{jmt}b_kX^{k}_{jmt}$ for $\ln b_k$, $-g_{jmt}E_{jmt}B_{jmt}$ for $\alpha$, $E_{jmt}B_{jmt}$ times the cell or market-quarter indicator for each effect, and $-\tilde{M}_{jmt}e^{\vartheta_{mt}}$ times the market-quarter indicator for each $\vartheta_{mt}$.

	Dividing every row by $E_{jmt}B_{jmt}>0$ leaves the column rank unchanged. The columns become $b_kX^{k}_{jmt}/B_{jmt}$, $-g_{jmt}$, the indicators of the effects, and $-e^{\vartheta_{mt}}\tilde{M}_{jmt}/(E_{jmt}B_{jmt})$ interacted with the market-quarter indicators. The effects' indicators span the bank $\times$ market and market $\times$ quarter means, so full column rank is the stated condition. The map from $b$ to $(\varsigma,\tilde{\kappa}_2,\tilde{\kappa}_1,\tilde{\kappa}_4,\varsigma_{\varnothing})$ is one to one.
\end{proof}

The condition has a plain reading. The ratio $X^{k}/B$ is the share of pool $k$ in the money that $j$'s rate pulls, so $b_k$ is read from how the net flow moves when that share changes within a cell over time. The level of $\tilde{\kappa}$ therefore rests on the within-cell variation of the pool weights. The market-quarter effect removes any movement common to the institutions of a market.

The checking scale and the checking pair enter only through $b_2$ and $b_4$. Own checking carries $e^{\varsigma}$ alone and rivals' checking carries $e^{\varsigma-\tilde{\kappa}_1}$, so the two are separated only by the variation of own against rival checking weights. The exit coefficient is separated from $\psi_{mt}$ because the exit scales with the own movable stock and the entry with every pool and with the price.

The proposition is about identification of $\tilde{\kappa}$, not of $\kappa$. What $\tilde{\kappa}$ measures depends on awareness and on how tastes sort across pools. A saver who prefers her institution is already there. If tastes for institutions are persistent and heterogeneous, savers at $j^\prime$ have on average a lower $\mu_{ij4mt}$ than savers at $j$ do, so they move to $j$ less for the same gap.

In \eqref{eq:20} that is a lower response of the rivals' pools for the same price, which raises $\tilde{\kappa}$ and looks like $\kappa$. It is the state dependence against unobserved heterogeneity problem of \textcite{heckman1981heterogeneity}, the same as in Panel B of Table \ref{tab:xi_diagnostics}. With $F_p$ the distribution of tastes among the holders of pool $p$, $F$ the market's, and $D(\mu_i)$ the common denominator at the depositor's tastes, sorting is
\begin{gather}
	h_{j\mid p}\equiv\frac{\int e^{\mu_{ij4mt}}D_p(\mu_i)^{-1}\,\text{d}F_p(d_{imt},\nu_{imt})}{\int e^{\mu_{ij4mt}}D(\mu_i)^{-1}\,\text{d}F(d_{imt},\nu_{imt})}, \label{eq:app_sorting}
\end{gather}
and $h_{jmt}$ is the pool-weighted value of $h_{j\mid p}$ over the rival pools of a pair, divided by its value for $j$'s own savings pool.

The share \eqref{eq:13} integrates over $F$ for every pool, so the model as written has $h\equiv1$, and sorting enters only if $F_p$ is tracked pool by pool. The fixed effects absorb the market-level mean of $\delta_{j4}$ but not the sorting of savers across institutions within a market.

\begin{prop}[What $\tilde{\kappa}$ measures]\label{prop:app_bound}
	Let \ref{assum:common_denominator} to \ref{assum:scale} hold, let tastes differ across pools as in \eqref{eq:app_Pexact} in place of \ref{assum:mean_utility}, and let savers at rival institutions consider $j$ with probability $A_{jmt}\in(0,1]$. Then the following hold.
	\begin{enumerate}[label=(\roman*)]
		\item The rival factor of the savings pair is $e^{-\tilde{\kappa}_2}$, with $\tilde{\kappa}_2=\kappa-\ln A_{jmt}-\ln h_{jmt}$. If awareness and sorting are the same for the daily-liquidity pair, $\tilde{\kappa}_4=\tilde{\kappa}_2-\ln E_c$.
		\item Under \ref{assum:sorting}, $\tilde{\kappa}_2\geq\kappa$.
		\item If $\kappa$ differs across pairs of institutions, (ii) holds with $\kappa$ replaced by $-\ln\sum_p w_pe^{-\kappa_p}$, where $w_p$ are the pool weights. That quantity is at most $\sum_p w_p\kappa_p$.
		\item Under \ref{assum:launch}, the change in $\tilde{\kappa}_2$ at 2020Q4 equals the change in $\kappa$.
	\end{enumerate}
\end{prop}

\begin{proof}
	(i) A saver in a rival pool $p$ chooses $(j,4)$ only if she considers $j$, so her probability is multiplied by $A_{jmt}$. Under \ref{assum:common_denominator}, the probability \eqref{eq:app_Pexact} of a rival pool $p$ is $c_pe^{\delta_{j4mt}-\kappa}$ times $\int e^{\mu_{ij4mt}}D(\mu_i)^{-1}\,\text{d}F_p$, which by \eqref{eq:app_sorting} is $h_{j\mid p}$ times the same integral over the market's $F$. That integral is common to every pool, own pools included, and goes into the entry scale. Relative to own savings, whose coefficient is one, the rival pools of the pair therefore enter with the factor $A_{jmt}h_{jmt}e^{-\kappa}$ in place of $e^{-\kappa}$.

	Under \ref{assum:market_size} the outside pool is measured, so its entry does not load on the rival weights. Under \ref{assum:scale} the scale is common to the own and the rival pools, so the ratio of their coefficients is the factor alone. For the daily-liquidity pair the clock multiplies the factor by $E_c$. (ii) $A_{jmt}\leq1$ and $h_{jmt}\leq1$ give $-\ln A_{jmt}\geq0$ and $-\ln h_{jmt}\geq0$.

	(iii) With costs that differ across pairs of institutions, the factor is $A_{jmt}h_{jmt}\sum_pw_pe^{-\kappa_p}$, and Jensen's inequality gives $\sum_pw_pe^{-\kappa_p}\geq e^{-\sum_pw_p\kappa_p}$. (iv) Differencing (i) across 2020Q4, the terms in $A_{jmt}$ and $h_{jmt}$ cancel when neither moves.
\end{proof}

Four remarks qualify the proposition. First, it keeps the common denominator and the common scale. In simulated data, random coefficients move the entry scale within cells, and the error of \ref{assum:common_denominator} pushes $\tilde{\kappa}$ down (Appendix \ref{app:simulation}). Where either happens, the bound may not hold.

Second, two mechanisms move $h$ in opposite directions. Holders who like their own institution more than the market does make $h<1$. Awareness that selects on taste, so that the savers who consider $j$ like it more, makes $h>1$ among those who consider $j$, which \ref{assum:sorting} rules out.

Third, if a share $\chi_{jmt}$ of rivals' money is held by savers who already have an account at $j$, the awareness term is $-\ln\bigl[\chi_{jmt}+A_{jmt}(1-\chi_{jmt})\bigr]$ rather than $-\ln A_{jmt}$. Fourth, a lower $\kappa$ after the launch may let savers re-sort, which would move $h$ and break \ref{assum:launch}. The change has not been simulated with random coefficients or with sorting, so the level is reported as a bound and the change rests on \ref{assum:launch}.

\subsection{Estimation and Inference}\label{app:netflow:estimation}

We estimate \eqref{eq:20} by nonlinear least squares, with each row divided by its total pool $1+Q^{1}+Q^{4}+R^{1}+R^{2}+R^{4}+\mathrm{Out}$. The weighting changes no parameter and is what makes the fit converge. Divided by own savings alone, the rows of institutions with almost no savings carry the whole sum of squares, in simulated and real data alike.

The effects are kept in levels and solved exactly by market, and the structural parameters by the method of Levenberg and Marquardt on the projected Jacobian, from four starts. The fit with the smallest sum of squares is kept, and where the starts end at an interior and at a boundary optimum both are reported.

The one endogenous regressor is $\rho_{j4mt}$ inside $g$, handled as Section \ref{sec:empirical:sleep} handles the spread, by a control function. Its first stage regresses the spread on the own cost shifters, and the residual enters the exponent.

Standard errors come from the wild cluster bootstrap by conglomerate \parencite{cameron2008bootstrap}, which carries the first stages of the control function and of the dose and holds $\omega_{jt}$, $\lambda_{jt}$ and the market size inside $\mathrm{Out}$ fixed. The savings pair has few destination conglomerates, so the bootstrap is the only inference to trust there.

Equation \eqref{eq:20} is estimated as it stands, nonlinear in the price. A linearized version fails in simulated data at every scale of the gaps, because the entry scale varies across cells (Appendix \ref{app:simulation}).

The sample is the bank $\times$ market $\times$ quarter cells of $\mathrm{B}$ institutions with a measured time-deposit rate at $t$ and $t-1$. Cells whose net flow, divided by own savings, exceeds its 99th percentile in absolute value are dropped, and the untrimmed sample is a robustness row. Cells whose rate is carried from an earlier quarter form a robustness row.

The maturity share $\omega_{jt}$ and the daily-liquidity share $\lambda_{jt}$ come from the institution's annual statements. The statement of year $y$ describes the stock carried into year $y+1$, so it applies to the quarters of $y+1$, and the quarters of 2016 take the 2016 statement. An institution without a statement takes the year's median among the $\mathrm{B}$ institutions with one, which is zero.

The first stage of the control function uses the conglomerate-quarters with a measured rate. In the Pix comparisons, the controls of the usage first stage, the Cad\'{U}nico share and log population, enter the second stage interacted with the launch in $\tilde{\kappa}$ and $\alpha$, so that they do not act as excluded instruments.

Advertising enters only in the quarters in which it is observed. A few observed values are negative. They are set to zero, and a robustness row drops them. The aggregation of the cells to bank $\times$ state and the national equation for $\mathrm{D}$ destinations are not estimated.

\subsection{The Pix Exposures}\label{app:netflow:pix}

We freeze the mandated set of \eqref{eq:21} at launch, so that exposure is fixed before Pix and does not respond to later adoption. The institutions that were not mandated are mid-size private and foreign banks, located mostly in large cities, and no state-owned bank is among them. Exposure therefore equals one in almost every market and varies only across a few large ones. Credit cooperatives do not report to ESTBAN, so the exposure is silent on the markets where they dominate.

The second exposure interacts the market's 2019 mobile access with the post-launch indicator. Mobile lines are already a market state in $\boldsymbol{S}_{mt}$, so this exposure identifies a Pix effect that varies with access. It is valid only if 2019 access is unrelated to other post-2020 shocks. Since the emergency aid went mostly to poorer markets, the design controls for Cad\'{U}nico registration interacted with the post-launch indicator.

\textcite{sarkisyan2024instant} relates the deposits of small relative to large banks in a municipality to local Pix use, and uses high-speed internet access, through 4G/5G mobile networks, as an alternative instrument. The third exposure follows him, with the same outcome and the 2019 4G/5G share, so that our estimates can be compared with his.

In the first stage of Table \ref{tab:pix_exposure_fs}, usage is the number of distinct individuals who send Pix transfers, per inhabitant, by MCA and quarter from 2020Q4 to 2024Q4. Each exposure is fixed at 2019 levels and scaled by its standard deviation. The regressions include quarter effects, 2019 Cad\'{U}nico registration interacted with quarter effects, and log population, with population weights.

\input{tab_pix_exposure_fs.tex} % Label: tab:pix_exposure_fs

\subsection{The Within-Bank Margin}\label{app:netflow:withinbank}

In the 2017 to 2018 Household Budget Survey (\textit{Pesquisa de Or\c{c}amentos Familiares}, POF), \pn[p0]{fees.pof.acct.share_paying} percent of families paid an account-keeping fee, R\$\pn[2]{fees.pof.acct.mean_payer} a month on average among those who paid, at January 2018 prices \parencite{ibge_pof2017_2019}. At the five largest banks (Banco do Brasil, Bradesco, Caixa, Ita\'{u}, and Santander), the median listed price of the first standardized service package rose from R\$\pn[2]{fees.listed.pkg1.big5.median.2016} in 2016 to R\$\pn[2]{fees.listed.pkg1.big5.median.2024} in 2024, in nominal terms \parencite{bcb_tarifas_mensal}. These are listed ceilings, not the prices clients pay.

On the public complaints platform \texttt{consumidor.gov.br}, fee disputes make up \pn[p0]{fees.disputes.acct_share.big5} percent of account complaints against the same five banks over 2016 to 2024, and \pn[p0]{fees.disputes.acct_share.digital} percent against Nubank, Inter, and C6 \parencite{senacon_consumidor_dados}. At the five banks the share falls to \pn[p0]{fees.disputes.acct_share.big5.2024} percent in 2024.

\section{Simulation Evidence}\label{app:simulation}

The simulation generates data from \eqref{eq:17} on the panel's own pools, with the mean utilities, the sleepy shares and the market sizes of Section \ref{sec:estimation} as the data-generating process, used for nothing else. Tastes are those of the logit, $\mu\equiv0$, except in the random-coefficients step below. The simulated world has no tax clock, so the time-deposit pair has the same true $\kappa$ as the other pairs.

The data carry no sampling noise and each design is simulated once, so the standard errors are cluster sandwiches of the misfit by conglomerate and a rough guide only. The switching cost is \pn[0]{kappa.k2b.truth.level} before 2020Q4 and changes by \pn[1]{kappa.k2b.truth.change} at the launch. In the dose version, $\kappa$ also takes \pn[2]{kappa.k2b.truth.dose} per unit of the market's Pix-use exposure after the launch. The price coefficient is that of Section \ref{sec:estimation:demand}, and it does not shift at the launch.

Two worlds are simulated. In the first, every institution offers time deposits and the common denominator of \ref{assum:common_denominator} holds exactly. In the second, an institution's time deposits count as a destination only where it holds them in the data, so the denominators differ across pools. Each fit uses the per-pair form of \eqref{eq:20}, with the rivals' time deposits in one pool, total-pool weighting, and two starts.

\begin{table}[H]
	\centering
	\begin{threeparttable}
		\caption{Recovery of the Switching Cost by Source Pair in Simulated Data}\label{tab:k2b_pairs}
		\input{tab_k2b_pairs.tex}
		\begin{tablenotes}[flushleft]
			\item \footnotesize \textit{Notes:} Estimates of \eqref{eq:20} by source pair, on data simulated from \eqref{eq:17} at the mean utilities, the sleepy shares and the market sizes of Section \ref{sec:estimation}. Level: the switching cost before the launch of Pix. Change: its shift at the launch. Dose: its coefficient on the market's Pix-use exposure after the launch. Ideal: every institution offers time deposits and the common denominator of Appendix \ref{app:netflow} holds exactly. FE: time deposits count as a destination only where the institution holds them in the data, as in the estimation panel with its fixed effects, so the denominators differ across pools. In the row labels, before-after is the fit with the launch only, and dose the fit that adds the exposure. Conv.: whether the fit converged. Standard errors in parentheses for the savings pair are cluster sandwiches by conglomerate. The simulated data carry no sampling noise, so they measure misfit and are a rough guide only. The FE dose fit stopped at the iteration cap because the checking pair ran to the edge of its range, printed as at bound. Its two starts agree on the savings pair.
		\end{tablenotes}
	\end{threeparttable}
\end{table}

Three results carry into the design. First, the savings pair recovers the change and the dose in both worlds. The change runs from \pn{kappa.k2b.sav.change.fe_ba} to \pn{kappa.k2b.sav.change.fe_dose} against \pn{kappa.k2b.truth.change}, and the dose is \pn{kappa.k2b.sav.dose.fe_dose} and \pn{kappa.k2b.sav.dose.ideal_dose} against \pn{kappa.k2b.truth.dose}. The level runs from \pn{kappa.k2b.sav.level.fe_dose} to \pn{kappa.k2b.sav.level.fe_ba} against \pn{kappa.k2b.truth.level}, below the truth by more than its standard error in three of the four fits, and the dose fit in the second world stopped at the iteration cap.

The checking pair misses even where \ref{assum:common_denominator} holds and runs to its bounds once the dose term is added, which is the weak identification of its scale rather than \ref{assum:common_denominator}'s error. The time-deposit pair's dose is \pn{kappa.k2b.td.dose.fe_dose} and \pn{kappa.k2b.td.dose.ideal_dose} against \pn{kappa.k2b.truth.dose}, on opposite sides of the truth.

Second, a linearized form of \eqref{eq:20} recovers nothing about $\kappa$ once the entry scale varies across cells, at every scale of the gaps tried, and dividing each row by its own savings pool instead of the total pool leaves the fit to a handful of rows and sends every parameter to its bound.

In a further world, in which the entry scale moves within a bank $\times$ market cell, the level is lost altogether. That world fails \ref{assum:scale} by construction and is not in the table. The change at the launch needs only that the scale does not shift within cells at 2020Q4 in a way correlated with the pools' composition, which was not simulated.

Third, the gross flows of the simulated world and of the data disagree. An awake R\$ stays with probability \pn[2]{kappa.k2.world.stay_prob.median} and awake money is \pn[p1]{kappa.k2.world.awake_share} percent of deposits in the simulation, against a net flow of \pn[p0]{kappa.k0.moving.median.k4} percent of the carried stock at the median time-deposit cell. The design is built on the net flow, which is observed.

Three breaks of the assumptions bear on the level, and Table \ref{tab:k2_breaks} reports them. The first replaces the logit tastes with the random coefficients of Section \ref{sec:estimation:demand}, with \pn{kappa.k2.rc.draws} draws per market, and fits one $\tilde{\kappa}$ shared by the three pairs on the second world's rows. With logit tastes these rows give \pn{kappa.k2.rc.k05.logit.kappa} and \pn{kappa.k2.rc.k10.logit.kappa} for a true $\kappa$ of \pn[1]{kappa.k2.rc.k05.truth} and \pn[0]{kappa.k2.rc.k10.truth}, below the truth, as a single $\tilde{\kappa}$ shared by the pairs is in this world.

With the random coefficients they give \pn{kappa.k2.rc.k05.rc.kappa} (\pn{kappa.k2.rc.k05.rc.se}) and \pn{kappa.k2.rc.k10.rc.kappa} (\pn{kappa.k2.rc.k10.rc.se}), neither fit converges, and the checking scale runs close to the edge of its range. The level fails under \ref{assum:mean_utility}. The random coefficients also move the entry scale within cells. The inflow they imply, relative to the logit inflow, is \pn[2]{kappa.k2.rc.k05.inflow_ratio.p10}, \pn[2]{kappa.k2.rc.k05.inflow_ratio.p50} and \pn[2]{kappa.k2.rc.k05.inflow_ratio.p90} at the 10th, 50th and 90th percentiles of cells, which breaks \ref{assum:scale}.

The second break sorts tastes toward the incumbent, with $h$ below one for the rival pools. The estimate then exceeds $\kappa$ in every case. With no switching cost it is close to $-\ln h$, and with a positive cost it exceeds $\kappa-\ln h$, by more the stronger the sorting and by up to \pn[2]{kappa.k2.sort.k05.max_excess}, so $\tilde{\kappa}-\kappa$ overstates the sorting term once the two combine.

The third break leaves the outside pool out of the equation. The estimate then falls to about zero whatever the true $\kappa$, since the rival savings pool takes the outside pool's entry. With the outside pool in, the savings pair gives \pn{kappa.k2.outside.k05.in.kappa} for a true $\kappa$ of \pn[1]{kappa.k2.rc.k05.truth} and \pn{kappa.k2.outside.k10.in.kappa} for \pn[0]{kappa.k2.rc.k10.truth}. At the largest cost it runs to the bound of its range, so that cell is not an estimate. A second start gives the same fits at the two largest costs, so these results reflect the form of the equation rather than the starting value. The second and third breaks use the per-pair form and the second world.

None of the three breaks was simulated for the change at the launch, which therefore rests on \ref{assum:launch} and on the simulations at the mean utilities above.

\begin{table}[H]
	\centering
	\begin{threeparttable}
		\caption{Breaking the Assumptions in Simulated Data}\label{tab:k2_breaks}
		\footnotesize
		\begin{tabular}{@{}l r r r r@{}}
			\toprule
			\multicolumn{5}{@{}l}{\textit{Panel A: random coefficients, one $\tilde{\kappa}$ shared by the pairs}} \\
			& \multicolumn{2}{c}{Logit tastes} & \multicolumn{2}{c}{Random coefficients} \\
			\cmidrule(lr){2-3}\cmidrule(lr){4-5}
			True $\kappa$ & $\tilde{\kappa}$ & Conv. & $\tilde{\kappa}$ & Conv. \\
			\pn[1]{kappa.k2.rc.k05.truth} & \pn{kappa.k2.rc.k05.logit.kappa} & \pn{kappa.k2.rc.k05.logit.converged} & \pn{kappa.k2.rc.k05.rc.kappa} & \pn{kappa.k2.rc.k05.rc.converged} \\
			 & (\pn{kappa.k2.rc.k05.logit.se}) & & (\pn{kappa.k2.rc.k05.rc.se}) & \\
			\pn[0]{kappa.k2.rc.k10.truth} & \pn{kappa.k2.rc.k10.logit.kappa} & \pn{kappa.k2.rc.k10.logit.converged} & \pn{kappa.k2.rc.k10.rc.kappa} & \pn{kappa.k2.rc.k10.rc.converged} \\
			 & (\pn{kappa.k2.rc.k10.logit.se}) & & (\pn{kappa.k2.rc.k10.rc.se}) & \\
			\midrule
			\multicolumn{5}{@{}l}{\textit{Panel B: tastes sorted toward the incumbent, savings pair}} \\
			True $\kappa$ & $h$ & $\kappa-\ln h$ & $\tilde{\kappa}$ & Conv. \\
			\pn[1]{kappa.k2.sort.k00.h10.truth} & \pn[1]{kappa.k2.sort.h10} & \pn{kappa.k2.sort.k00.h10.truth} & at zero & \pn{kappa.k2.sort.k00.h10.converged} \\
			 & \pn[1]{kappa.k2.sort.h09} & \pn{kappa.k2.sort.k00.h09.truth} & \pn{kappa.k2.sort.k00.h09.kappa} (\pn{kappa.k2.sort.k00.h09.se}) & \pn{kappa.k2.sort.k00.h09.converged} \\
			 & \pn[1]{kappa.k2.sort.h07} & \pn{kappa.k2.sort.k00.h07.truth} & \pn{kappa.k2.sort.k00.h07.kappa} (\pn{kappa.k2.sort.k00.h07.se}) & \pn{kappa.k2.sort.k00.h07.converged} \\
			 & \pn[1]{kappa.k2.sort.h05} & \pn{kappa.k2.sort.k00.h05.truth} & \pn{kappa.k2.sort.k00.h05.kappa} (\pn{kappa.k2.sort.k00.h05.se}) & \pn{kappa.k2.sort.k00.h05.converged} \\
			\pn[1]{kappa.k2.sort.k05.h10.truth} & \pn[1]{kappa.k2.sort.h10} & \pn{kappa.k2.sort.k05.h10.truth} & \pn{kappa.k2.sort.k05.h10.kappa} (\pn{kappa.k2.sort.k05.h10.se}) & \pn{kappa.k2.sort.k05.h10.converged} \\
			 & \pn[1]{kappa.k2.sort.h09} & \pn{kappa.k2.sort.k05.h09.truth} & \pn{kappa.k2.sort.k05.h09.kappa} (\pn{kappa.k2.sort.k05.h09.se}) & \pn{kappa.k2.sort.k05.h09.converged} \\
			 & \pn[1]{kappa.k2.sort.h07} & \pn{kappa.k2.sort.k05.h07.truth} & \pn{kappa.k2.sort.k05.h07.kappa} (\pn{kappa.k2.sort.k05.h07.se}) & \pn{kappa.k2.sort.k05.h07.converged} \\
			 & \pn[1]{kappa.k2.sort.h05} & \pn{kappa.k2.sort.k05.h05.truth} & \pn{kappa.k2.sort.k05.h05.kappa} (\pn{kappa.k2.sort.k05.h05.se}) & \pn{kappa.k2.sort.k05.h05.converged} \\
			\midrule
			\multicolumn{5}{@{}l}{\textit{Panel C: the outside pool left out, savings pair}} \\
			& \multicolumn{2}{c}{Outside pool in} & \multicolumn{2}{c}{Outside pool out} \\
			\cmidrule(lr){2-3}\cmidrule(lr){4-5}
			True $\kappa$ & $\tilde{\kappa}$ & Conv. & $\tilde{\kappa}$ & Conv. \\
			\pn[1]{kappa.k2.sort.k00.h10.truth} & at zero & \pn{kappa.k2.outside.k00.in.converged} & at zero & \pn{kappa.k2.outside.k00.out.converged} \\
			\pn[1]{kappa.k2.rc.k05.truth} & \pn{kappa.k2.outside.k05.in.kappa} (\pn{kappa.k2.outside.k05.in.se}) & \pn{kappa.k2.outside.k05.in.converged} & \pn{kappa.k2.outside.k05.out.kappa} (\pn{kappa.k2.outside.k05.out.se}) & \pn{kappa.k2.outside.k05.out.converged} \\
			\pn[0]{kappa.k2.rc.k10.truth} & \pn{kappa.k2.outside.k10.in.kappa} (\pn{kappa.k2.outside.k10.in.se}) & \pn{kappa.k2.outside.k10.in.converged} & \pn{kappa.k2.outside.k10.out.kappa} (\pn{kappa.k2.outside.k10.out.se}) & \pn{kappa.k2.outside.k10.out.converged} \\
			\pn[0]{kappa.k2.outside.k20.truth} & box & \pn{kappa.k2.outside.k20.in.converged} & \pn{kappa.k2.outside.k20.out.kappa} (\pn{kappa.k2.outside.k20.out.se}) & \pn{kappa.k2.outside.k20.out.converged} \\
			\bottomrule
		\end{tabular}
		\begin{tablenotes}[flushleft]
			\item \footnotesize \textit{Notes:} Estimates of \eqref{eq:20} on data simulated from \eqref{eq:17} in the world where time deposits count as a destination only where the institution holds them (FE in Table \ref{tab:k2b_pairs}). Panel A: the data-generating process uses the random coefficients of Section \ref{sec:estimation:demand} with \pn{kappa.k2.rc.draws} draws per market, or logit tastes on the same rows, and the fit has one $\tilde{\kappa}$ for the three source pairs. The fits use \pn{kappa.k2.rc.n} rows in \pn{kappa.k2.rc.markets} markets. Conv.: one if the fit converged, zero if it stopped at the iteration cap. Panel B: tastes sorted toward the incumbent, so that the estimate should equal $\kappa-\ln h$. Panel C: the outside pool in and out of the equation. Panels B and C use the per-pair form, with standard errors in parentheses beside the estimate. In Panel C a second start reproduces the fits at the two largest values of $\kappa$. At zero: the fit sits at the lower bound of $\tilde{\kappa}$. Box: the fit runs to the upper bound of $\tilde{\kappa}$, so the cell is not an estimate. Standard errors in parentheses are cluster sandwiches by conglomerate and measure misfit, since the simulated data carry no sampling noise.
		\end{tablenotes}
	\end{threeparttable}
\end{table}

\end{document}

% ============================================================================================================================
% Bibliography: four new entries for References.bib. Web sources, so no Google Scholar BibTeX exists for the first three.
% ============================================================================================================================
% @misc{bcb_pix_participantes,
%   author       = {{Banco Central do Brasil}},
%   shortauthor  = {BCB},
%   title        = {Lista de Participantes do {Pix}},
%   date         = {2020/2024},
%   howpublished = {Banco Central do Brasil},
%   url          = {https://www.bcb.gov.br/content/estabilidadefinanceira/pix/ListadeparticipantesdoPix.pdf},
%   urldate      = {2026-06-09},
%   annotation   = {Dated lists, 22 October 2020 to 11 November 2024, each marking an institution's participation as mandatory (Resolu{\c{c}}{\~a}o BCB n{\textordmasculine} 1, 2020) or voluntary. All lists retrieved from the Internet Archive, captures of 2020 to 2026.}
% }
% @misc{bcb_pix_chaves,
%   author       = {{Banco Central do Brasil}},
%   shortauthor  = {BCB},
%   title        = {Pix: Estoque de Chaves por Participante ({ChavesPix}), Dados Abertos},
%   date         = {2020/2024},
%   howpublished = {Banco Central do Brasil, Olinda},
%   url          = {https://olinda.bcb.gov.br/olinda/servico/Pix_DadosAbertos/versao/v1/odata/ChavesPix(Data=@Data)},
%   urldate      = {2026-10-27},
%   annotation   = {Month-end stock of Pix keys by participant (ISPB), person type and key type.}
% }
% @misc{datafolha_topofmind,
%   author       = {{Folha de S.Paulo}},
%   title        = {Top of Mind: Finan{\c{c}}as},
%   date         = {2016/2022},
%   howpublished = {Folha de S.Paulo, survey by Datafolha},
%   url          = {https://top-of-mind.folha.uol.com.br/},
%   urldate      = {2026-11-05},
%   annotation   = {Annual national survey of spontaneous brand recall; bank category, editions of October 2016 to October 2022. Editions from 2023 are behind the paywall.}
% }
% @article{rothenberg1971identification,
%   title   = {Identification in parametric models},
%   author  = {Rothenberg, Thomas J},
%   journal = {Econometrica},
%   volume  = {39},
%   number  = {3},
%   pages   = {577--591},
%   year    = {1971}
% }
